% !TEX root = ../../../main.tex
% 来源：中学数学实验教材·第二册（几何）
% 章节：实验几何 / 全等和叠合
% 原始文件：1.tex，行 2460-2517
% 类型：tikz
% ---
\begin{tikzpicture}
\begin{scope}
\tkzDefPoints{0/0/A, 2/0/B, 1.5/1/C}
\tkzDrawPolygon(A,B,C)
\node at (1.3,.5){I};
\end{scope}
\begin{scope}[xshift=2.3cm]
	\tkzDefPoints{0/0/A, 2/0/B, 1.3/1/C}
\tkzDrawPolygon(A,B,C)
\node at (1.3,.5){II};
\end{scope}
\begin{scope}[scale=1.3, xshift=3.5cm]
	\tkzDefPoints{0/0/A, 2/0/B, 1.5/1/C}
\tkzDrawPolygon(A,B,C)
\node at (1.2,.5){III};
\node at (1,-.5){(1)};
\end{scope}
\begin{scope}[xshift=7.5cm]
	\tkzDefPoints{0/0/A, 1.5/0/B, 1/1.5/C}
\tkzDrawPolygon(A,B,C)
\node at (.8,.7){IV};
\end{scope}
\begin{scope}[xshift=9.5cm]
	\tkzDefPoints{0/0/A, 2/0/B, 1.5/1/C}
\tkzDrawPolygon(A,B,C)
\node at (1.3,.5){V};
\end{scope}
\begin{scope}[yshift=-2.5cm]
	\tkzDefPoints{0/0/A, 2/0/B, 1.3/1.2/C, 1.6/-.5/D}
	\tkzDrawPolygon(A,D,B,C)
	\tkzDrawSegments[dashed](A,B)
	\tkzDrawSegments(C,D)
\node at (1.1,.5){$\rm I'$};
\end{scope}
\begin{scope}[yshift=-2.5cm,xshift=3cm]
	\tkzDefPoints{0/0/A, 2/0/B, 1.1/1.1/C, 1.3/-.6/D}
	\tkzDrawPolygon(A,D,B,C)
	\tkzDrawSegments[dashed](A,B)
	\tkzDrawSegments(C,D)
	\node at (1.1,.5){$\rm II'$};
	\node at (3,-1){(2)};
\end{scope}
\begin{scope}[yshift=-2.5cm,xshift=6cm]
	\tkzDefPoints{0/0/A, 2/0/B, 0.7/1.1/C, 1.3/-.6/D}
	\tkzDrawPolygon(A,D,B,C)
	\tkzDrawSegments[dashed](A,B)
	\tkzDrawSegments(C,D)
	\node at (.9,.5){$\rm III'$};
\end{scope}
\begin{scope}[yshift=-2.5cm,xshift=9cm]
	\tkzDefPoints{0/0/A, 2/0/B, 1.1/1.1/C, 1.3/-.6/D}
	\tkzDrawPolygon(A,D,B,C)
	\tkzDrawSegments[dashed](A,B)
	\tkzDrawSegments(C,D)
	\node at (1.1,.5){$\rm IV'$};
\end{scope}

\end{tikzpicture}
